% !TEX root = main.tex
\documentclass[main.tex]{subfiles}
\begin{document}

\section{Incompressible Lagrangian flows}
\label{sec:flow}

We now turn the constrained power diagrams of \S\ref{sec:contributions} into a Lagrangian discretization of incompressible flows, in which each cell is a fluid parcel of fixed mass moving with velocity $\vel_i \!\in\! \R^d$.
For simplicity, we assume a uniform unit density, so the mass of each cell equals its prescribed volume $\pvol_i$, and we collect the masses into the diagonal matrix $\smash{\pvolmat = \mathrm{diag}(\pvol_i \Id_d)}$.
We denote the time step by $\tstep$.
We first review Power Particles \cite{deGoes2015PowerParticles}, which serves as our baseline (\S\ref{sec:flow-pp}), and the penalty formulation of \citet{Gallouet2018Lagrangian} (\S\ref{sec:flow-penalty}), which both take the sites as particles.
We then derive a variational formulation that takes the centroids as particles, and whose time steps are exactly the softened problem of \S\ref{sec:regularized-joint} (\S\ref{sec:flow-centroids}), before relating its Newton solver to the operators of Power Particles (\S\ref{sec:flow-newton}).

\subsection{Baseline: Power Particles}
\label{sec:flow-pp}

Power Particles \cite{deGoes2015PowerParticles} store, for each particle, a site $\qq_i$ and a velocity $\vel_i$, and advance them by an operator splitting.
Given sites $\qq^k$, volumes $\pvol$, and velocities $\vel^k$, each time step performs:
\begin{enumerate}
  \item \emph{Volume enforcement:} compute the weights $\weight^k$ that maximize $\dualenergy$ in \eqref{eq:volume-dual} for the sites $\qq^k$, so that the power cells have the prescribed volumes $\pvol$, and evaluate their centroids $\centroid^k$.
  \item \emph{Pressure projection:} solve the Poisson equation $\smash{\tstep\,\Lmat\,\press = \Dmat\,\vel^k}$ and update $\smash{\vel^{k+1} = \vel^k - \tstep\,\pvolmat^{-1} \Gmat\,\press}$, where $\Dmat$ is the discrete divergence of the cell volumes with respect to the sites at fixed weights, $\smash{\Gmat = -\Dmat\tr}$ its discrete gradient, and $\smash{\Lmat = \Dmat\,\pvolmat^{-1} \Gmat}$ the pressure Laplacian \cite[Eqs.~(4) and~(7)]{deGoes2015PowerParticles}.
  \item \emph{Advection:} displace the sites using the centroids and the projected velocities, $\smash{\qq^{k+1} = \centroid^k + \tstep\,\vel^{k+1}}$.
\end{enumerate}
External forces are added to the velocities before the projection, and the listing above is a cyclic rotation of \cite[Alg.~1]{deGoes2015PowerParticles}, which solves for the weights at the end of each step.
Incompressibility is enforced twice: dynamically, since the projected velocity is discretely divergence-free, $\smash{\Dmat\,\vel^{k+1} = 0}$, so the volumes are preserved to first order when the sites move at fixed weights; and geometrically, since the weight solve restores the exact volumes after every advection.
Advecting the sites from the centroids rather than from the current sites keeps the cells well shaped, and it drives particles at rest towards a centroidal power diagram.
However, this splitting is not derived from a variational principle, and, as noted by \citet[Remark~1.5]{Gallouet2018Lagrangian}, it lacks a Hamiltonian structure and a convergence analysis, although it does not depend on a penalty parameter.
The other formulations below are compared against this baseline.

\subsection{Sites as particles: a penalty formulation}
\label{sec:flow-penalty}

\citet{Gallouet2018Lagrangian} proposed a Lagrangian scheme in which the particles are again the sites $\qq$, and the cells are volume-constrained, \ie the power diagram given by the optimal weights of \eqref{eq:volume-dual}.
Incompressibility is relaxed into a stiff penalty on the squared Wasserstein distance between the uniform measure on $\domain$ and the particles, which, for fixed volumes $\pvol_i$, equals the transport cost $\smash{\transport(\qq) = \max_\weight \dualenergy(\weight; \qq)}$.
The resulting Lagrangian reads
\begin{equation}
  \lag(\qq, \dot\qq)
    = \frac{1}{2} \sum_i \pvol_i \norm{\dot\qq_i}^2
    - \frac{1}{2\stiff^2}\, \transport(\qq) ,
  \label{eq:gm-lagrangian}
\end{equation}
where $\stiff \!>\! 0$ controls the stiffness of the penalty.
By the envelope theorem and \eqref{eq:centroid-gradient}, the site gradient of the transport cost is $\smash{2\,\pvol_i (\qq_i - \centroid_i)}$, where $\centroid_i$ is the centroid of the volume-constrained cell $\cell_i$.
The equations of motion thus pull each site toward its own centroid \cite[Eq.~(1.7)]{Gallouet2018Lagrangian},
\begin{equation}
  \ddot\qq_i = -\frac{1}{\stiff^2} \bigl( \qq_i - \centroid_i \bigr) ,
  \label{eq:gm-ode}
\end{equation}
so a centroidal power diagram with prescribed volumes is an equilibrium of the flow.
Integrating \eqref{eq:gm-ode} with a symplectic Euler step yields \cite[Eq.~(1.11)]{Gallouet2018Lagrangian}
\begin{subequations}
  \label{eq:gm-step}
  \begin{empheq}[left={\empheqlbrace\,}]{align}
    \vel^{k+1}_i &= \vel^k_i + \frac{\tstep}{\stiff^2} \bigl( \centroid^k_i - \qq^k_i \bigr) , \label{eq:gm-step-vel} \\[2mm]
    \qq^{k+1}_i &= \qq^k_i + \tstep\, \vel^{k+1}_i . \label{eq:gm-step-pos}
  \end{empheq}
\end{subequations}
Their convergence analysis towards smooth solutions of the Euler equations assumes a time step no larger than the stiffness, $\tstep \!\le\! \stiff$, and its error bound vanishes only for $\tstep \!=\! o(\stiff^2)$ \cite[Theorem~1.4 and Remark~1.6]{Gallouet2018Lagrangian}.

The two site-based schemes are closely related.
Substituting \eqref{eq:gm-step-vel} into \eqref{eq:gm-step-pos} with the stiffness $\stiff \!=\! \tstep$ gives
\begin{equation}
  \qq^{k+1}_i = \centroid^k_i + \tstep\, \vel^k_i ,
  \label{eq:gm-advection}
\end{equation}
which is the advection step of Power Particles, except that it uses the velocity \emph{before} the update.
In this light, the site-to-centroid gap $\smash{(\centroid^k_i - \qq^k_i)/\tstep}$ in \eqref{eq:gm-step-vel} plays the role of a pressure impulse.
Power Particles discards this impulse from the velocity and replaces it by the pressure projection, which makes the velocities exactly divergence-free, while \citet{Gallouet2018Lagrangian} kept it as the only pressure force, which makes the velocities divergence-free only in the stiff limit $\stiff \!\to\! 0$.
Note, however, that the choice $\stiff \!=\! \tstep$ lies at the boundary of the regime covered by their convergence analysis, and outside the regime where their error bound vanishes.

\subsection{Centroids as particles: a variational formulation}
\label{sec:flow-centroids}

Both schemes above evolve the sites, while the centroids are only a by-product of the volume-constrained power diagram.
We now take the opposite view: the configuration is a partition $\cell$ of $\domain$ with the prescribed volumes $\pvol$, and the particles are the centroids $\centroid$ of its cells, which carry the masses.

\paragraph{Potential.}
For any sites $\qq$ and any partition with volumes $\pvol$, the transport cost splits as
\begin{equation}
  \sum_i \int_{\cell_i} \norm{\xx - \qq_i}^2 d\xx = \sum_i \int_{\cell_i} \norm{\xx - \centroid_i}^2 d\xx + \sum_i \pvol_i \norm{\centroid_i - \qq_i}^2 .
  \label{eq:flow-split}
\end{equation}
The penalty of \citet{Gallouet2018Lagrangian} minimizes the left-hand side over partitions for sites placed at the particles.
With the particles at the centroids instead, the last term vanishes, and we take as potential the remaining term,
\begin{equation}
  \cpot(\cell) = \frac{1}{2\stiff^2} \sum_i \int_{\cell_i} \norm{\xx - \centroid_i}^2 d\xx ,
  \label{eq:flow-cpot}
\end{equation}
which is the quantization energy of the partition.
Note that the transport cost $\transport$ itself cannot serve as the potential of the centroids: $\transport(\centroid)$ is attained by the optimal partition for the sites $\centroid$, whose centroids differ from $\centroid$, so the particles would no longer be the centroids of their cells, as in \S\ref{sec:flow-penalty}.
Since reassigning the cells to the sites $\centroid$ by optimal transport and then recomputing their centroids can only decrease $\cpot$, its minimizers are centroidal power diagrams with the volumes $\pvol$, which are the configurations favored by the centroid-based advection of Power Particles.

\paragraph{Action and time steps.}
The motion of the particles follows from the action of the Lagrangian
\begin{equation}
  \lag = \frac{1}{2} \sum_i \pvol_i \norm{\dot\centroid_i}^2 - \cpot(\cell) .
  \label{eq:flow-lagrangian}
\end{equation}
To discretize the action in time, we first advect the particles with their current velocities to the positions
\begin{equation}
  \qq^k_i = \centroid^k_i + \tstep\,\vel^k_i ,
  \label{eq:flow-advect}
\end{equation}
where external forces $\mathbf{f}_i$ are included by adding $\smash{\tstep^2\,\mathbf{f}_i / \pvol_i}$.
We then compute the next partition by minimizing the incremental potential of an implicit Euler step over the partitions with volumes $\pvol$,
\begin{equation}
  \cell^{k+1} = \argmin_{\cell} \ \frac{1}{2\tstep^2} \sum_i \pvol_i \norm{\centroid_i - \qq^k_i}^2 + \cpot(\cell) ,
  \label{eq:flow-implicit}
\end{equation}
and update the velocities by finite differences of the centroids, $\smash{\vel^{k+1} = (\centroid^{k+1} - \centroid^k) / \tstep}$.

\paragraph{Equivalence with softened centroid constraints.}
Applying the splitting \eqref{eq:flow-split} with the sites $\qq^k$, the objective of \eqref{eq:flow-implicit} scaled by $2\stiff^2$ becomes
\begin{equation*}
  \sum_i \int_{\cell_i} \norm{\xx - \qq^k_i}^2 d\xx + \alpha \sum_i \pvol_i \norm{\centroid_i - \qq^k_i}^2 ,
\end{equation*}
with the stiffness
\begin{equation}
  \alpha = \frac{\stiff^2}{\tstep^2} - 1 .
  \label{eq:flow-alpha}
\end{equation}
This is exactly the softened problem \eqref{eq:regularized-primal} of \S\ref{sec:regularized-joint}, with both the sites and the targets set to the advected positions $\qq^k$.
Its first term is the transport cost of \citet{Gallouet2018Lagrangian} evaluated at the advected positions, so each time step minimizes this transport cost over the partitions while penalizing the deviation of their centroids from the advected positions.
For $\tstep \!\le\! \stiff$, the stiffness is non-negative, which is the regime of the convergence analysis of \citet{Gallouet2018Lagrangian}, and each time step amounts to maximizing the concave dual $\smash{\jointdual_\alpha(\shift, \weight; \qq^k, \qq^k)}$ of \eqref{eq:regularized-dual}.
At the maximizer, the cells form a power diagram with the exact volumes $\pvol$, generated by the shifted sites $\smash{\dual = \qq^k + \shift}$, and \eqref{eq:regularized-gradient-shift} gives $\smash{\shift = \alpha\,(\qq^k - \centroid^{k+1})}$.
The velocity update then reads
\begin{equation}
  \vel^{k+1} = \vel^k - \frac{1}{\alpha\,\tstep}\, \shift ,
  \label{eq:flow-velocity}
\end{equation}
so the shifts are the pressure impulses of the particles.
Equivalently, since $\smash{-\shift = \qq^k - \dual}$, the impulse on each particle points from the generator $\dual_i$ of its cell toward its advected position $\qq^k_i$, just as the impulse \eqref{eq:gm-step-vel} of \citet{Gallouet2018Lagrangian} points from the generator of each cell toward its centroid.
The cells form a power diagram as long as the generators remain pairwise distinct, which the condition \eqref{eq:regularized-distinct} guarantees when $\stiff$ is close enough to $\tstep$ and the advected positions lie in the convex hull of $\domain$; beyond this regime, distinct generators are not guaranteed.

\paragraph{Algorithm.}
In summary, given the centroids $\centroid^k$ and velocities $\vel^k$, each time step performs:
\begin{enumerate}
  \item \emph{Advection:} compute the sites $\qq^k$ by \eqref{eq:flow-advect}, as in step~3 of Power Particles.
  \item \emph{Softened solve:} maximize $\smash{\jointdual_\alpha(\shift, \weight; \qq^k, \qq^k)}$, which yields cells with the volumes $\pvol$, their centroids $\centroid^{k+1}$, and the shifts $\shift$.
  \item \emph{Velocity update:} compute $\vel^{k+1}$ by \eqref{eq:flow-velocity}.
\end{enumerate}
The sites are thus not independent variables: they are recomputed at every step from the particles and their velocities by the advection rule of Power Particles, and the generators of the diagram are these sites displaced by the shifts.

\paragraph{Relation to the site-based schemes.}
The stiffness \eqref{eq:flow-alpha} relates our scheme to those of \S\ref{sec:flow-pp} and \S\ref{sec:flow-penalty}.
In the limit $\alpha \!\to\! 0$, \ie $\stiff \!\to\! \tstep$, the shifts vanish and $\jointdual_\alpha$ reduces to the volume-constrained dual $\dualenergy(\weight; \qq^k)$ (\S\ref{sec:regularized-joint}).
The new centroids are then the centroids $\smash{\hat\centroid}$ of the volume-constrained power diagram of the advected sites, as computed by the weight solve of Power Particles (step~1), and, since $\smash{\shift / \alpha = \qq^k - \centroid^{k+1}}$, the velocity update becomes
\begin{equation*}
  \vel^{k+1} = \vel^k + \frac{1}{\tstep} \bigl( \hat\centroid - \qq^k \bigr) ,
\end{equation*}
which is the impulse \eqref{eq:gm-step-vel} of \citet{Gallouet2018Lagrangian} with $\stiff \!=\! \tstep$.
The limit $\alpha \!\to\! 0$ is thus Power Particles with its pressure projection (step~2) replaced by the impulse of \citet{Gallouet2018Lagrangian}.
For $\stiff \!>\! \tstep$, the shifts no longer vanish, and they correct the centroids so as to balance the kinetic term against the potential, while keeping the volumes exact.

\paragraph{Semi-implicit variant.}
For partitions with volumes $\pvol$, the splitting \eqref{eq:flow-split} with the sites $\centroid^k$ shows that $\smash{\tfrac{1}{2\stiff^2} \sum_i \int_{\cell_i} \norm{\xx - \centroid^k_i}^2 d\xx}$ is exactly the linearization of the potential $\cpot$ at the current centroids $\centroid^k$.
Treating the potential explicitly thus gives the time step
\begin{equation}
  \cell^{k+1} = \argmin_{\cell} \ \frac{1}{2\tstep^2} \sum_i \pvol_i \norm{\centroid_i - \qq^k_i}^2 + \frac{1}{2\stiff^2} \sum_i \int_{\cell_i} \norm{\xx - \centroid^k_i}^2 d\xx ,
  \label{eq:flow-semi-implicit}
\end{equation}
which is convex for any time step.
This step is again the softened problem of \S\ref{sec:regularized-joint}, now with \emph{distinct} sites and targets: the sites are the current particles $\centroid^k$, the targets are the advected positions $\qq^k$, and the stiffness is $\smash{\alpha = \stiff^2 / \tstep^2}$, \ie each step maximizes $\smash{\jointdual_\alpha(\shift, \weight; \centroid^k, \qq^k)}$.
Each step thus minimizes the transport cost to the current particles while softly matching their advected positions.
Moreover, whenever the previous cells form a power diagram, the sites $\centroid^k$ are automatically pairwise distinct, since two convex cells with disjoint interiors are separated by a hyperplane and their centroids lie strictly on either side.

\paragraph{Weakly compressible flows.}
Softening the volume constraints as in \S\ref{sec:softened-volumes} turns the time step into a weakly compressible one, in which each weight follows the equation of state \eqref{eq:softened-volume-state} with stiffness $\beta$, reminiscent of the compressible variant of Power Particles \cite[Eq.~(9)]{deGoes2015PowerParticles}.
The identity between the kinetic term and the centroid penalty then breaks, since the slack constraint of \eqref{eq:regularized-primal-cons} gives $\smash{\pvol_i \norm{\slack_i}^2 = (\cellvol_i^2 / \pvol_i) \norm{\centroid_i - \tcentroid_i}^2}$, which overweighs compressed cells.
This mismatch can be removed by penalizing instead the \emph{perspective} $\smash{\norm{\mathbf{m}_i}^2 / \cellvol_i = \cellvol_i \norm{\centroid_i - \tcentroid_i}^2}$ of the centroid moment $\smash{\mathbf{m}_i = \int_{\cell_i} (\xx - \tcentroid_i)\,d\xx}$, which is jointly convex in the moments and volumes, and thus in the (fractional) partition.
Its conjugate representation $\smash{\alpha \norm{\mathbf{m}_i}^2 / \cellvol_i = \max_{\shift_i} -2\,\shift_i\tr \mathbf{m}_i - \tfrac{1}{\alpha} \cellvol_i \norm{\shift_i}^2}$ is linear in the partition, and its last term only adds $\smash{\norm{\shift_i}^2 / \alpha}$ to the combined weight of \eqref{eq:shifted-site-weight}.
The partition minimizer is therefore again a power diagram, and the dual objective becomes
\begin{equation}
  \jointdual(\shift, \weight + \tfrac{1}{\alpha} \norm{\shift}^2) - \frac{1}{\alpha} \sum_i \pvol_i \norm{\shift_i}^2 - \frac{1}{4\beta} \sum_i \pvol_i\,\weight_i^2 ,
  \label{eq:flow-perspective-dual}
\end{equation}
where $\smash{\norm{\shift}^2}$ denotes the vector of squared shift norms; it is concave as a minimum of concave functions of the multipliers, and its shift gradient is $\smash{2\,\cellvol_i (\tcentroid_i - \centroid_i) - \tfrac{2}{\alpha}\,\cellvol_i\,\shift_i}$, so that $\smash{\shift_i = \alpha\,(\tcentroid_i - \centroid_i)}$ holds regardless of the volume error.
For exact volumes, \eqref{eq:flow-perspective-dual} is merely the softened dual $\jointdual_\alpha$ with the weights reparametrized by $\smash{\weight \mapsto \weight + \norm{\shift}^2/\alpha}$, so both formulations share their maximizing shifts and partition.
With the perspective penalty, the time step weighs both the kinetic term and the transport cost, whose centroid part is $\smash{\sum_i \cellvol_i \norm{\centroid_i - \qq^k_i}^2}$, by the current volumes, \ie it treats the masses as proportional to the current cell volumes.
Weighing the kinetic term by the fixed masses instead, $\smash{\pvol_i \norm{\centroid_i - \qq^k_i}^2 = \pvol_i \norm{\mathbf{m}_i}^2 / \cellvol_i^2}$, would make the primal non-convex in the partition, so the perspective penalty is the closest convex alternative, with a first-order relative error of $\smash{(\cellvol_i - \pvol_i)/\pvol_i}$, half of that of \eqref{eq:regularized-primal}.

\paragraph{Reference weights.}
The softened dual \eqref{eq:softened-volume-dual} penalizes the weights about zero, so each time step starts its pressure from scratch.
Given reference weights $\bweight$, \eg the weights of the previous time step, we can instead write the weights as $\smash{\weight = \bweight + \delta\weight}$ and penalize only the deviation $\delta\weight$, which turns the softened dual into
\begin{equation*}
  \jointdual_{\alpha}(\shift, \weight) - \frac{1}{4\beta} \sum_i \pvol_i\,(\weight_i - \bweight_i)^2 ,
\end{equation*}
which is the dual of the softened primal whose transport cost uses the power distances $\smash{\norm{\xx - \qq_i}^2 - \bweight_i}$ of the weighted sites $(\qq, \bweight)$ in place of \eqref{eq:regularized-primal-obj}.
The equation of state \eqref{eq:softened-volume-state} then holds for the deviation, $\smash{\delta\weight_i = 2\beta\,(\pvol_i - \cellvol_i) / \pvol_i}$, so the pressure accumulates across time steps instead of being recomputed from the current volume error alone.
Adding a common constant to $\bweight$ leaves the cells unchanged, since $\smash{\sum_i \pvol_i = \domainVol}$, and in the limit $\beta \!\to\! 0$ the weights are pinned to $\bweight$, which recovers $\centroiddual_\alpha$ of \eqref{eq:softened-centroid-dual} with the fixed weights $\bweight$.
The same change of variables applies to \eqref{eq:flow-perspective-dual}.

\subsection{Newton solver and pressure projection}
\label{sec:flow-newton}

We compute each time step \eqref{eq:flow-implicit} by Newton's method on $\jointdual_\alpha$, and its linear system turns out to assemble the operators of Power Particles.
In this section, we drop the time index and write $\qq$ for the advected sites $\qq^k$.
We initialize the shifts at zero and the weights at the maximizer of $\dualenergy(\weight; \qq)$, which is the weight solve of Power Particles after advection.
We write the negated Hessian of $\jointdual_\alpha$ in blocks, using the derivatives of \S\ref{sec:joint-constraints} and \S\ref{sec:regularized-joint}.
The shift--shift block is $\smash{\Kmat + \tfrac{2}{\alpha}\,\pvolmat}$, with $\smash{\Kmat = -\nabla^2_{\shift\shift} \centroiddual}$ the negated centroid Hessian \eqref{eq:centroid-hessian}.
The weight--weight block $\smash{\wlap = -\nabla^2_{\weight\weight} \jointdual}$ is the negated volume Laplacian \eqref{eq:newton-step}.
The mixed block is the discrete divergence $\smash{\Dmat = \partial \cellvol / \partial \shift}$ of \S\ref{sec:joint-constraints}.
Each Newton step thus solves
\begin{equation}
  \begin{bmatrix} \Kmat + \tfrac{2}{\alpha}\,\pvolmat & \Dmat\tr \\ \Dmat & \wlap \end{bmatrix}
  \begin{bmatrix} \delta\shift \\ \delta\weight \end{bmatrix}
  =
  \begin{bmatrix} \residual \\ \pvol - \cellvol \end{bmatrix} ,
  \label{eq:flow-newton}
\end{equation}
whose right-hand side is the gradient of $\jointdual_\alpha$, with $\smash{\residual_i = 2\,\cellvol_i (\qq_i - \centroid_i) - \tfrac{2}{\alpha}\,\pvol_i\,\shift_i}$ by \eqref{eq:regularized-gradient-shift}.
This matrix is positive definite for $\alpha \!>\! 0$ up to the additive constant on the weights, since the shift block is positive definite and $\wlap$ is a Laplacian.

The blocks have a direct interpretation at the initial shifts $\shift \!=\! 0$.
There, the shifted sites are the advected sites, $\dual \!=\! \qq$, and by \eqref{eq:appendix-jac} the combined weights do not vary to first order with the shifts, since $\smash{\partial_{\shift_i} \combweight_i = 2(\dual_i - \tcentroid_i) = 0}$.
A shift thus acts as a displacement of a site at fixed weights, so $\Dmat$ is \emph{exactly} the discrete divergence of Power Particles \cite[Eq.~(4)]{deGoes2015PowerParticles}, evaluated on the power diagram of the advected sites $\qq$, and $\smash{\Dmat\tr = -\Gmat}$ is its discrete gradient.
Likewise, $\wlap$ is the finite-volume Laplacian that Power Particles uses for its weight solve \cite[Eq.~(3)]{deGoes2015PowerParticles}, up to the normalization of the weights.
If the diagram is moreover close to centroidal, the shift block is approximately $\smash{\Kmat \approx 2\,\pvolmat\,\Jmat}$, where $\smash{\Jmat = \partial \centroid / \partial \dual}$ is the Jacobian of the cell centroids with respect to the sites.

Eliminating the shifts, the first row of \eqref{eq:flow-newton} gives
\begin{equation}
  \delta\shift = \bigl( \Kmat + \tfrac{2}{\alpha}\,\pvolmat \bigr)^{-1} \bigl( \residual + \Gmat\,\delta\weight \bigr) ,
  \label{eq:flow-newton-shift}
\end{equation}
and the second row yields the Schur complement system for the weights,
\begin{equation}
  \Bigl( \wlap - \Dmat \bigl( \Kmat + \tfrac{2}{\alpha}\,\pvolmat \bigr)^{-1} \Dmat\tr \Bigr) \delta\weight
    = \bigl( \pvol - \cellvol \bigr) - \Dmat \bigl( \Kmat + \tfrac{2}{\alpha}\,\pvolmat \bigr)^{-1} \residual .
  \label{eq:flow-newton-schur}
\end{equation}
To relate these equations to Power Particles, assume that $\Jmat$ is close to the identity on modes that are smooth at the scale of the cells, as it is for uniform translations.
Then $\smash{\Kmat + \tfrac{2}{\alpha}\,\pvolmat \approx 2\dilation\,\pvolmat}$ with $\smash{\dilation = 1 + 1/\alpha}$, and, at the initial shifts, where the volumes are exact and $\smash{\residual = 2\,\pvolmat (\qq - \hat\centroid)}$ with the centroids $\smash{\hat\centroid}$ of the volume-constrained power diagram of $\qq$, the first Newton step yields the velocity update \eqref{eq:flow-velocity}
\begin{equation}
  \vel^{k+1} \approx \vel^k + \frac{\tstep}{\stiff^2} \bigl( \hat\centroid - \qq \bigr) - \tstep\,\pvolmat^{-1} \Gmat\,\press ,
  \qquad
  \press = \frac{1}{2\stiff^2}\,\delta\weight ,
  \label{eq:flow-newton-velocity}
\end{equation}
using $\smash{\alpha\,\dilation = \stiff^2 / \tstep^2}$.
The first correction is the impulse \eqref{eq:gm-step-vel} of \citet{Gallouet2018Lagrangian}, evaluated at the advected sites $\qq$, and the second is a discrete pressure gradient, as in the velocity correction of Power Particles \cite[Alg.~1, line~3]{deGoes2015PowerParticles}, with the weight update in the role of the pressure.
The pressure solves \eqref{eq:flow-newton-schur}, whose matrix becomes
\begin{equation}
  \wlap + \frac{1}{2\dilation}\,\Lmat , \qquad \Lmat = \Dmat\,\pvolmat^{-1} \Gmat ,
  \label{eq:flow-newton-laplacians}
\end{equation}
where $\Lmat$ is the pressure Laplacian of Power Particles \cite[Eq.~(7)]{deGoes2015PowerParticles}, and whose right-hand side is the divergence $\smash{\tfrac{1}{\dilation}\,\Dmat \bigl( \hat\centroid - \qq \bigr)}$ of the impulse.
Without the weight Laplacian $\wlap$, this system would read $\smash{\tstep\,\Lmat\,\press = \Dmat \bigl( \tfrac{\tstep}{\stiff^2} (\hat\centroid - \qq) \bigr)}$, \ie the pressure projection of Power Particles applied to the impulse, in which case the velocity change in \eqref{eq:flow-newton-velocity} would be discretely divergence-free.
Since $\wlap$ is positive semi-definite while $\Lmat$ is negative semi-definite, the matrix \eqref{eq:flow-newton-laplacians} is the difference of the one-ring and two-ring Laplacians; to leading order on smooth modes, $\smash{\wlap \approx -\tfrac12 \pvol \nabla^2}$ and $\smash{\Lmat \approx \pvol \nabla^2}$, so their difference approximates $\smash{-\tfrac{1}{2(1+\alpha)} \pvol \nabla^2}$, whose smooth part weakens as the centroids are pinned more strongly, $\alpha \!\to\! \infty$, leaving mostly oscillatory modes; we have not quantified this approximation.
The weight Laplacian accounts for the freedom to reshape the cells while keeping their volumes, which Power Particles exploits only in its separate weight solve.

To first order, the time step therefore reorders the splitting of \S\ref{sec:flow-pp}: it first advects the sites with the unprojected velocity and enforces their volumes, as in the steps~3 and~1 of Power Particles, and then projects the resulting impulse with a pressure that combines the one-ring Laplacian $\wlap$ and the two-ring Laplacian $\Lmat$, in place of step~2.
As $\alpha \!\to\! 0$, the shift updates vanish, the pressure term drops, and only the impulse remains, as discussed in \S\ref{sec:flow-centroids}.
Unlike Power Particles, which linearizes the incompressibility constraint at the current diagram and then performs a separate nonlinear weight solve, Newton's method iterates on a single strongly concave problem, and the identities above hold exactly only at its first step, since subsequent steps evaluate the operators at shifted sites.
We have not quantified the approximation $\smash{\Jmat \approx \Id}$, which neglects the deviation of the cells from a centroidal configuration.

\end{document}
